\chapter{The conjecture, special cases, and status}
\label{ch:conjecture-and-examples}

\section{The property for one group}
\begin{definition}\label{def:baum-connes-property}
\uses{def:assembly-map}
\notready
Write $\mathrm{BC}(\Gamma)$ for the proposition that $\mu_{\Gamma,i}$
is an isomorphism for both $i=0$ and $i=1$.
\end{definition}

\section{The general conjecture}
\begin{conjecture}\label{conj:baum-connes}
\uses{def:baum-connes-property}
\notready
Every countable discrete group $\Gamma$ satisfies $\mathrm{BC}(\Gamma)$.
\end{conjecture}

This is a target statement, not an axiom or an available theorem.
Its eventual Lean statement must use the actual constructed assembly map.

\section{Calibration examples and later specializations}
Plan separate examples for the trivial group, finite groups, and $\mathbb Z$.
Treat $\mathrm{SL}(3,\mathbb Z)$ as a later specialization of the general
framework. For each example separately track constructions, K-group
computations, identification of assembly, injectivity, and surjectivity.

\section{Research and formalization status}
Keep established mathematics, completed Lean work, and open mathematical
questions separate. This structural reorganization provides no new
mathematical result or proof mechanism for the general conjecture.
